\subsection{Molecular Geometries and Active Space Preparation}
The dataset used comprised 28 closed-shell organic molecules from Thiel's excited-state benchmark set \cite{thiels2008benchmark}. The published cartesian nuclear coordinates were referenced directly and used without further optimisation. 
For each molecule, Restricted Hartree-Fock calculations were performed using PySCF \cite{pyscf2018,pyscf2020} with the STO-3G basis set. This set of canonical molecular orbitals were optimised once and used consistently throughout all subsequent CASCI, VQE and QSE calculations in the study. 

Active spaces (denoted $n_e e\, n_{\mathrm{o}} o$) comprised of the $n_e / 2$ occupied molecular spatial-orbitals closest to the HOMO and the $n_{\mathrm{o}} - n_e / 2$ virtual molecular spatial-orbitals cloest to the LUMO. The remaining occupied orbitals were treated as frozen and doubly occupied, while the remaining virtual orbitals were excluded from calculations. 

\subsection{Exact Reference}
Within each active space, the exact reference energies and wavefunctions were obtained from complete active-space configuration interaction (CASCI) calculations performed using PySCF. 
Calculations were separated by spin-projection sector $M_S$ by setting the number of active $\alpha$ and $\beta$ spin electrons to be 

$n_\alpha = \frac{n_e}{2} + M_S, n_\beta = \frac{n_e}{2} - M_S$

Singlet states were obtained from the $M_S=0$ sector, while triplet states were calculated in the $M_S=0, +1, -1$ sectors. For each active-space spin projection sector, the determinant-space dimension, and therefore the maximum number of available roots, was $\binom{n_{\mathrm{o}}}{n_\alpha} \binom{n_{\mathrm{o}}}{n_\beta}$. 

Where available, the lowest 11 singlet states (ground state + 10 excited states), and the lowest 10 triplet states (within each spin projection sector) were calculated. To ensure that the target number of states were obtained, three times the required number of CASCI roots were initialy requested from the FCI solver, subject to the maximum dimension of the corresponding determinant space. 
The resulting roots were then filtered according to their $\Ssq$ expectation values, retaining only roots that satisfy $\left| \Ssq - S(S+1) \right| < 10^{-5}$ with $S=0$ for singlets and $S=1$ for triplets. 

Throughout this study, wavefunctions were represented in a common active spin-orbital determinant basis by a state vector of length $\binom{2 n_{\mathrm{o}}}{n_{\mathrm{e}}}$. 
This vector includes all determinants with $n_e$ electrons distributed amoung the $2n_{\mathrm{o}}$ active spin-orbitals (irespective of their spin projections).
Determinants were ordered lexicographically using an interleaved spin-orbital convention $(1\alpha, 1\beta, 2\alpha, 2\beta...)$.

PySCF represented each wavefunction as a two dimensional CI coefficient matrix with dimensions $\binom{n_{\mathrm{o}}}{n_{\mathrm{\alpha}}}$ by $\binom{n_{\mathrm{o}}}{n_{\mathrm{\beta}}}$, coressponding to the $\alpha$ and $\beta$ electron occupation strings. These coefficients were embedded into the corresponding entries of the common state-vector representation.  During this conversion, the fermionic phase arising from the permutation between PySCF's separated $\alpha$ and $\beta$ electron ordering and the interleaved spin-orbital convention was included explicitly.


\subsection{Quantum Computing Framework}
All quantum-circuit calculations in this study were performed using classical simulations implemented with the Qibo and Qibochem \cite{qibo2021,qibochem2024zenodo} frameworks. These frameworks were used to construct, represent and simulate the quantum circuits employed.
A research fork of Qibochem was developed to support the computaional workflow used in the study. The fork introduced configurable ansatz classes, optimised circuit evaluation routines, support for measuring spin observables, and dedicated functionality for constructing and sampling the projected matrices required for Quantum Subspace Expansion calculations. 

All calculations in this work used ideal noiseless circuit simulations without modelling any hardware. Only finite-sampling noise arising from a limited number of measurement shots was considered in the relevant sections.

\subsubsection{Hamiltonian and Spin Operators}
For each active space, the electronic Hamiltonian was constructed in second quantisation as follows:
\ElectronicHamiltonianEquation
where $E_{\mathrm{nuc}}$ is the core nuclear repulsion energy and $p, q, r, s$ index the active spin orbitals. $h_{pq}$ and $g_{pqrs}$ are the one and two-electron integral coefficients evaluated in the fixed canonical Hartree-Fock MO basis.

Spin operators were constructed in the same second quantisation basis as follows.

\SpinZEquation
\SpinLadderOperatorsEquation
\SpinSquaredEquation

All fermionic Hamiltonian, spin and excitation operators were mapped to qubit operators using the Jordan-Wigner (JW) transformation.

\subsubsection{Ground State Reference}
As introduced above, approximate ground-state reference wavefunctions were prepared using parameterised Unitary Coupled Cluster (UCC) Ansatz. 
In this study, 6 classes of Ansatz were used, namely: UCCSD, UCCSDSinglet, UCCGSD, and \(k\)-UpCCGSDSinglet with \(k=1,2,\) and \(3\). The Ansatz were characterised by the excitation generators included in the cluster operator $\hat{T}$, summarised in the following table.

\input{tables/methods/ansatz}

For each ansatz, candidate singles and doubles transitions were first enumerated in fermionic form. Only canonical transitions that excite from a lower to high spin orbital were retained because the reverse transitions would already be included through its conjugate adjoint $\hat{T}^{\dagger}$ before exponentiation. All retained generators were spin conserving: single excitations preserved spin labels, while double excitations preserved the total number of $\alpha$ and $\beta$ spin electrons. 

The non-generalised UCCSD and UCCSDSinglet ansatz were restricted to occupied-to-virtual single and double excitations relative to the reference Hartree-Fock determinant \cite{Anand2022UCCReview}. In contrast, the generalised UCCGSD ansatz allowed generalised orbital transitions between occupied--occupied and virtual--virtual spin orbitals \cite{Lee2018GeneralisedAnsatzOriginal, Diniz2021GeneralisedAnsatzAdaptations}.
To reduce the size of the generalised-double operator pool, mix-direction double excitaitons, in which one electron was promoted while the other was relaxed, were excluded. Across the $4e4o$ to $8e8o$ active spaces, this practical restriction reduced the number of independent parameters by 42.2--45.2\% and the circuit depth by 46.4--47.6\%.

The $k$-UpCCGSDSinglet ansatz combined generalised singles with generalised paired doubles. For each paired doubles, an $\alpha\beta$ electron pair was removed from one spatial orbital and created in another. 
The UCCSDSinglet and $k$-UpCCGSDSinglet constructions imposed spin-adapted parameter sharing between spin-orbital excitations that had the same spatial-orbital labels. 
For $k$-UpCCGSDSinglet, the complete generalised single and paired-double block was repeatd $k$ times, with an independent set of variational paramters assigned to each repetition. 

Unless otherwise stated, all UCC operators were implemented using the first order Suzuki Trotter product approximation given in \cref{eq:first-order-suzuki-trotter}. 

\SuzukiTrotterFirstOrderEquation

At every optimisation step, the circuit energy was evaluated as the noiseless statevector expectation value of the electronic Hamiltonian \cref{eq:electronic-hamiltonian}
The variational parameters were optimised classically using the L-BFGS-B optimiser through Qibo, to minimise the ground state expectation value. For the $k$-UpCCGSDSinglet ansatz, the parameters associated with all $k$ repeated layers were optimised simultaneously.
Before optimisation, ansatz parameters were initialised based on the following hierarchy: 
\[
\mathrm{MP2}
\rightarrow
1\text{-UpCCGSD}
\rightarrow
2\text{-UpCCGSD}
\rightarrow
3\text{-UpCCGSD}
\rightarrow
\mathrm{UCCSDSinglet}
\rightarrow
\mathrm{UCCSD}
\rightarrow
\mathrm{UCCGSD}.
\]

Parameters associated with excitation generators shared with the preceding ansatz were initialised with the previously optimised amplitudes, while newly introduced generators were initialised to zero. The initial occupied to virtual double-excitation parameters were estimated using second-order Moller-Plesset perturbation theory (MP2). For parameter transfer from a \(k\)-UpCCGSDSinglet ansatz to a subsequent ansatz, the optimized amplitudes associated with each generator were summed across the \(k\) layers and used as the initial amplitude for the corresponding generator

Following optimisation, the resulting VQE parameters were stored as a fixed ground-state reference wavefunction. This fixed reference set was used consistently for all subsequent statevector and finite-shot QSE calculations.

\subsubsection{QSE Operator Pools and Projected Matrices}
As discussed in the introduction, separate one-body operator pools were used to construct the singlet and triplet-targetting QSE subspaces. The singlet operator pool, $\mathcal{O}^{S}$, was defined as

\QSESingletOperatorPoolEquation
wherbey, $i$ and $j$ to denote active spatial orbitals.

The triplet operator pool, $\mathcal{O}^{T}$, was constructed from all three spin-projection components, which were collected into a single operator pool.
\QSETripletOperatorPoolEquation

The singlet operator pool therefore contained $n_{\mathrm{o}}^2$
operators, while the triplet pool contains
$3n_{\mathrm{o}}^2$ operators. Individual pool elements were enumerated as
$\{\hat{O}_\mu\}$, with $\mu$ and $\nu$ indexing arbitrary operators within the given pool.

A projected observable matrix was first constructed for the overlap, hamiltonian and spin-squared matrices. Observable matrix elements were defined as follows: 
\QSEProjectedObservablesEquation
where the corresponding target observables were $\hat I, \hat H,$ and $\hat S^{2}$, respectively.

The fermionic QSE expansion operators and each target observable were mapped separately to Pauli operators using the Jordan-Wigner transformation. The projected observables were then constructed by multiplication in the Pauli representation. Terms with coefficients of magnitude smaller than $10^{-12}$ were discarded after the final product construction. 
The resulting Pauli-string expectation values were then evaluated with respect to the optimised VQE states, $\left|\Psi_{\mathrm{VQE}}\right\rangle$
For all calculations performed, only the diagonal and upper-triangular matrix elements, $\mu <= \nu$ were constructed and evaluated. Lower-triangular matrix terms were then reconstructed using Hermitian symmetry.


\subsubsection{Finite-Shot Sampling}
Finite-shot sampling was used to model the effects of statistical measurement noise on the projected Hamiltonian and overlap matrices. 
For each calculation, all unique Pauli strings in the projected Hamiltonian and Overlap observable matrices were first collected into a global measurement pool. Then, a largest-first greedy qubit-wise-commuting (QWC) grouping algorithm was used to partition the global measurement pool. 
Grouping was performed globally across elements from both matrices which reduced the total number of commuting groups that needed to be measured. Preliminary data also showed that this method of global grouping, as opposed to grouping terms for each matrix element, resulted in much better agreement with the statevector results. 

For each commuting group, a separate measurement circuit was constructed by appending basis-rotation gates to the optimized VQE circuit, such that all Pauli strings in the group could be measured in the computational ($Z$) basis. Hadamard gates were applied to qubits measured in the ($X$) basis, while $S^{\dagger}$ followed by a Hadamard gate was applied to qubits measured in the ($Y$) basis. This produced a distinct rotated statevector for each commuting group.
A sampled probability vector was sampled for every rotated statevector using Numpy's multinomial sampler. The expectation values of ever Pauli string in each group were then calculated based on the sampled probability vector. $H$ and $S$ matrix elements were then re-constructed from the expectations of every Pauli string in the global measurement pool. 

Measurement outcomes were generated for each rotated statevector using NumPy's multinomial sampler. The expectation values of all Pauli strings within a commuting group were estimated from the resulting sampled frequencies. These expectation values were then used to reconstruct the Hamiltonian and overlap matrices.

At each specified shot count, the finite-shot calculation was repeated using ten statistically independent sampling iterations. The resulting Hamiltonian and overlap matrices were solved independently, and the reported energy errors were calculated from the mean of the ten independently obtained energies.

\subsection{Numerical Analysis and Evaluation Metrics}
\subsubsection{Generalised Eigenvalue Problem}

To obtain the energy eigenvalues and wavefunctions from the QSE matrices, we must solve the following equation Hc=ESc


For every eigenvalue 

Solving the problem 

Define thresholding and dimension solve 

stable dimensions (2e3o etc) -> Why need to filter even for SV 

Change of basis to pvec form for projection calculations 

\subsubsection{Spin-expectation Evaluation}

\subsubsection{Projection}
Formula to calculate projection 

\subsubsection{Triplet Degenerecies}
Grouping degenerate states 
Mean of projection 
Splitting 

\subsubsection{Non Index State Matching}
hungarian algorithm 
